\chapter{Einführung}
\label{c:intro}
Durch die steigende Popularität von Chatbots wie ChatGPT von OpenAI und Claude von Anthropic in 
den letzten Jahren, wie man in \ref*{fig:trend_gpt} sehen kann, werden die Fortschritte, die im 
Bereich der künstlichen Intelligenz erreicht werden, immer bekannter und der Nutzen den KI 
Applikationen im Leben von Menschen haben wird dadurch immer erweitert.\par
Chatbots übernehmen inzwischen auch Aufgaben im Bereich der Kommunikation 
mit Kunden. Manche Unternehmen haben statt direktem Kundensupport einen 
Support Bot, über den Kunden Informationen erhalten können oder den 
Kundenservice kontaktieren können, falls es nicht möglich ist das Problem
nur durch den Bot zu lösen.\par


\section{Motivation}
\label{s:intro:mot}
Die Motivation für diese Bachelorarbeit entspringt der Erkenntnis, dass
trotz konstanter Forstschritte im Bereich der Large Language Models und 
Chatbots trotzdem weiterer Forschungs- und Entwicklungsbedarf besteht.\par

Das Thema dieser Arbeit wurde zusammen mit der Techdivision GmbH erarbeitet.
Die Techdivision hat interne Support Charts, für die Mitarbeiter der Firma, 
in denen Sie bei technischen Problemen Fragen stellen können um vom 
IT-Support Team Hilfe erhalten. Einige dieser Fragen sind Routinefragen, 
welche regelmäßig gestellt werden und teils mithilfe der internen Dokumentation
gelöst werden könnten. Um die Ressourcen des IT-Teams zu schonen und 
den Mitarbeitern bei einer schnelleren Lösung ihres Problems zu helfen,
soll ein Chatbot erstellt werden, der mithilfe interner Daten Fragen 
beantwortet und somit das IT-Support Team unterstützt. \par
Diese Arbeit zielt darauf ab, einen Beitrag zur Überbrückung der Lücke
 zwischen theoretischen Konzepten und praktischer Anwendung zu leisten. 
 Durch die konzeptionelle Erstellung eines Chatbots und die anschließende 
 prototypische Implementierung wird ein Ansatz verfolgt, 
 der es ermöglicht, verschiedene LLMs im 
 Kontext realer Anwendungsszenarien zu evaluieren.\par
Der Vergleich von vier unterschiedlichen LLMs innerhalb des entwickelten 
Prototyps bietet die Gelegenheit, die Leistungsfähigkeit, 
Effizienz und Anwendbarkeit dieser Modelle zu untersuchen. Diese 
vergleichende Analyse soll nicht nur zur Ermittlung des am besten 
geeigneten Modells für diesen spezifischen Anwendungsfall beitragen, sondern 
auch wertvolle Einblicke in die generellen Charakteristika und 
Einsatzmöglichkeiten verschiedener LLMs liefern.\par
Darüber hinaus adressiert diese Arbeit die wachsende Nachfrage nach 
praxisorientierten Lösungen im Bereich der KI-gestützten Kommunikation. 
In einer Zeit, in der Unternehmen und Organisationen zunehmend auf KI-Systeme 
setzen, um Effizienz zu steigern und Kundenerlebnisse zu verbessern, 
liefert diese Arbeit konkrete Erkenntnisse für die 
Entwicklung und den Einsatz moderner Chatbot-Systeme.\par
\iffalse
\begin{figure}[htbp]
    \centering
    \begin{tikzpicture}
        \begin{axis}[
            width=\textwidth,
            height=0.6\textwidth,
            xlabel={Week},
            ylabel={ChatGPT Interest},
            xlabel near ticks,
            ylabel near ticks,
            date coordinates in=x,
            xticklabel style={
                rotate=45,
                anchor=north east,
                inner sep=0pt,
                outer sep=2pt,
            },
            date ZERO=2022-06-26,
            xmajorgrids=true,
            ymajorgrids=true,
            grid style={dotted,gray!30},
            legend pos=north west,
            title={ChatGPT Google Trends},
        ]
            \addplot[
                thick,
                blue,
                mark=*,
                mark size=1pt,
            ] table [x=Woche, y=ChatGPT, col sep=comma] {data/ChatGPT.csv};
            \addlegendentry{ChatGPT}
            \addplot[
                thick,
                red,
                mark=*,
                mark size=1pt,
            ] table [x=Woche, y=Claude, col sep=comma] {data/Claude.csv};
            \addlegendentry{Claude}
        \end{axis}
    \end{tikzpicture}
    \caption{ChatGPT und Claude Interesse (Google Trends Data vom 26.06.2022 bis 26.05.2024)}
    \label{fig:trend_gpt}
\end{figure}
\fi
\iffalse
\begin{figure}[htbp]
    \centering
    \begin{tikzpicture}
        \begin{axis}[
            width=\textwidth,
            height=0.6\textwidth,
            xlabel={Week},
            ylabel={Claude Interest},
            xlabel near ticks,
            ylabel near ticks,
            date coordinates in=x,
            xticklabel style={
                rotate=45,
                anchor=north east,
                inner sep=0pt,
                outer sep=2pt,
            },
            date ZERO=2022-06-26,
            xmajorgrids=true,
            ymajorgrids=true,
            grid style={dotted,gray!30},
            legend pos=north west,
            title={ChatGPT Google Trends},
        ]
            \addplot[
                thick,
                blue,
                mark=*,
                mark size=1pt,
            ] table [x=Woche, y=Claude, col sep=comma] {data/ChatGPT.csv};
            \addlegendentry{Claude}
        \end{axis}
    \end{tikzpicture}
    \caption{ChatGPT Interest Over Time (Google Trends Data)}
    \label{fig:trend_gpt}
\end{figure}
\fi
\iffalse
Durch den Aufstieg von verschiedenen textbasierten Chat-Modellen wie OpenAIs Chat-GPT oder Googles Gemini, hat man in 
den letzten Jahren gesehen, zu was Modelle inzwischen in der Lage sind. Sie sind nicht mehr nur in der Lage, mit Menschen 
zu kommunizieren, sondern helfen diesen auch bei ihren Jobs oder Schulangelegenheiten. Verschiedene trainierte Modelle 
sind inzwischen in der Lage, komplizierte Probleme zu lösen und Menschen das Leben einfacher zu machen. 

Die Techdivision ist intern in sogenannten Services strukturiert, die häufig ein intensives Zusammenarbeiten erfordern. 
Diese Zusammenarbeit findet über verschiedene Kommunikationskanäle statt, darunter JIRA Boards und Support-Chats. Oft 
werden Routinefragen gestellt, die automatisiert werden könnten. 

Dadurch stellt sich speziell in den Support-Chats die Frage, wie viele der Supportanfragen sich durch einen KI-gesteuerten 
Chatbot ohne Hilfe eines Mitarbeiters lösen lassen würden und inwiefern ein Chatbot zur Entlastung der Teams und zur 
schnelleren Lösung von Problemen beitragen könnte.
Durch die Automatisierung einfacher Supportanfragen könnte zum einen Zeit gespart, zum Anderen auch die Ressourcen der 
Teams effektiver genutzt werden. Mitarbeiter könnten sich verstärkt auf die schwierigeren Supportanfragen konzentrieren, 
während der Chatbot die repetitiven und einfachen klärt, was nicht nur das Team entlastet, sondern auch die Reaktionszeit 
auf Supportanfragen erhöht.
\fi
\section{Zielsetzung}
\label{s:intro:goal}
\iffalse
Das Ziel der Arbeit besteht darin, einen Chatbot zu konzipieren und einen Prototyp zu entwickeln, der mit Hilfe der Daten 
der internen Datenquellen der Techdivision Support leistet und Routinefragen beantwortet. Der Chatbot soll idealerweise 
innerhalb von Google Chat integriert werden. Dabei ist es wichtig, dass der Chatbot, statt in Einzelchats, in einer Chat 
Gruppe läuft. Diese beinhaltet neben den Support Suchenden, ebenfalls Mitarbeiter des jeweiligen Services. Auf diese Weise 
kann, falls notwendig, extern Eingegriffen werden. Außerdem sollte der Chatbot in der Lage sein zu überprüfen, wie viele 
der aktuellen Supportanfragen er eigenständig beantworten oder lösen kann. Falls der Chatbot keine Antwort auf eine Anfrage 
haben sollte, soll er in den angebundenen Datenquellen nachschlagen können, um mit dem Wissen seine Wissensbasis zu 
erweitern und die Fragen entsprechend zu beantworten.
\fi
Das Ziel der Arbeit besteht darin, einen Chatbot zu konzipieren und einen Prototyp zu entwickeln, der mit Hilfe der Daten 
der internen Datenquellen der Techdivision Support leistet und Fragen beantwortet. Der Chatbot soll 
innerhalb von Google Chat integriert werden. Dabei ist es wichtig, dass der Chatbot, statt in Einzelchats, in einer Chat 
Gruppe läuft. Diese beinhaltet neben den Support Suchenden, ebenfalls Mitarbeiter des jeweiligen Services. Auf diese Weise 
kann, falls, notwendig, extern eingegriffen werden. Außerdem sollte der Chatbot in der Lage sein zu überprüfen, wie viele 
der aktuellen Supportanfragen er eigenständig beantworten oder lösen kann. Falls der Chatbot keine Antwort auf eine Anfrage 
haben sollte, soll er in den angebundenen Datenquellen nachschlagen können, um mit dem Wissen seine Wissensbasis zu 
erweitern und die Fragen entsprechend zu beantworten.\par
Hierfür wird im Rahmen dieser Arbeit ein Konzept für einen Chatbot erstellt, 
welches prototypisch implementiert wird. Dazu werden über die Implementierung vier verschiedene LLMs 
verwendet und miteinander verglichen, um das beste für diesen Anwendungsfall zu ermitteln. 
Der Vergleich wird in zwei Stufen durchgeführt. In der ersten Stufe werden 
die LLMs über das Framework Ragas verglichen. Im zweiten Schritt werden
die zwei besten LLMs aus dem Ragas Vergleich, über einen 
Fragebogen Vergleich von Mitarbeitern der TD bewertet und miteinander verglichen.